\exo
On considère la fonction $g$ définie sur $\mathbb{R}^*$ par $g(x)=\dfrac{1}{x}$.
À l'aide de sa représentation graphique, donner une valeur approchée des images suivantes.
\setlength{\columnseprule}{0pt}
\begin{multicols}{2}
\begin{enumerate}
\item $g(4)$
\item $g(0,25)$
\item $g(-3,2)$
\item $g\left(\dfrac{3}{5}\right)$
\item $g\left(-\dfrac{7}{5}\right)$
\end{enumerate}
\end{multicols}

Déterminer ensuite, graphiquement, un antécédent lorsqu'il existe.
\begin{multicols}{3}
\begin{enumerate}
\item $g(x)=4$
\item $g(x)=-0,44$
\item $g(x)=\dfrac{2}{5}$
\end{enumerate}
\end{multicols}

\exo
Résoudre graphiquement dans $\mathbb{R}^*$ les inéquations suivantes.
\begin{multicols}{2}
\begin{enumerate}
\item $\dfrac{1}{x}\leq -2$
\item $1\leq \dfrac{1}{x}\leq 2$
\item $0\leq \dfrac{1}{x}\leq 3$
\item $-2\leq \dfrac{1}{x}\leq 0$
\item $4>\dfrac{1}{x}\geq -1$
\end{enumerate}
\end{multicols}

\exo
Résoudre dans $\mathbb{R}^*$ les équations suivantes.
\begin{multicols}{2}
\begin{enumerate}
\item $\dfrac{1}{x}=2$
\item $\dfrac{1}{x}=-5$
\item $\dfrac{3}{x}=2$
\item $\dfrac{1}{x}=\dfrac{5}{3}$
\item $\dfrac{1}{x}+3=2$
\item $\dfrac{2}{x}=x$
\item $\dfrac{7}{x}+2=5$
\end{enumerate}
\end{multicols}

\exo
Pour chacune des fonctions suivantes, déterminer le sens de variation sur chacun des intervalles de son ensemble de définition, puis dresser le tableau de variations.
\begin{enumerate}
\item $f$ définie sur $\mathbb{R}^*$ par $f(x)=\dfrac{1}{x}+3$.
\item $g$ définie sur $\mathbb{R}^*$ par $g(x)=-\dfrac{2}{x}$.
\end{enumerate}
